\section{Introduction}

In recent years, large language models (LLMs) have grown substantially in both popularity and capabilities, becoming increasingly important tools in everyday life. Modern LLMs are capable of following complex user instructions with little to no additional guidance~\cite{brown2020fewshotlearners, wei2022finetuned} and have achieved significant improvements across a wide range of tasks, including programming~\cite{Chen2021EvaluatingLL}, mathematics and reasoning~\cite{mathds, cot}, question answering~\cite{brown2020fewshotlearners}, and natural language understanding~\cite{bert}.

While LLMs demonstrate remarkable performance across many areas, they are also prone to generating content that is considered harmful or undesirable, such as vulnerable code~\cite{pearce2022copilot} or toxic content~\cite{gehman2020realtoxicityprompts, bot}. To mitigate the generation of such responses, LLMs typically undergo safety alignment during post-training~\cite{tulu3, ouyang2022rlhf, rafailov2023dpo}, with the objective of aligning model outputs with human preferences and safety objectives.

While safety alignment has generally made LLMs more likely to produce harmless responses, the safety-aligned models have been shown to remain vulnerable to jailbreak attacks~\cite{dan2024, zou2023GCG, liu2024autodan, andriushchenko2025simpleadaptive, ding2023renellm}, where attackers circumvent safety-aligned behavior and induce the model to generate harmful responses to malicious prompts that it would otherwise refuse.
Recent works have proposed a wide range of jailbreak attacks, including methods that use imaginary scenarios~\cite{dan2024, dualuse, chao2025pair}, obfuscation~\cite{dualuse}, role-playing~\cite{johnny2024}, and other prompt modifications~\cite{xie2025sorrybench} to conceal or disguise the malicious intent of the original prompt.

Moreover, more recently, as open-weight models have become increasingly common, more powerful attacks have been proposed that expose the fragility of safety alignment in LLMs. Multiple recent works~\cite{wu2026neurostrike, twinbreak, wei2024brittleness} have shown that safety alignment in LLMs is highly sparse, with only a few thousand neurons, commonly referred to as safety-critical neurons, playing a central role in safety-aligned behavior. These works have shown that pruning these neurons can effectively remove safety alignment and induce harmful responses while having negligible impact on the model's general capabilities.
Another recent study~\cite{arditi2024refusal} has shown that safety-aligned refusal behavior in large language models can be effectively disabled by ablating a single direction from the model's intermediate activations and representations.%, which, we hypothesize, can be conceptually attributed to the sparse representation of safety alignment.

Taken together, these results suggest that current safety alignment depends on a surprisingly small subset of neurons, creating a structural weakness that is fundamentally different from conventional prompt-level jailbreak vulnerabilities. These findings are particularly concerning for open-weight models, where attackers have unrestricted access to model parameters and internal representations and can directly target the components responsible for safety-aligned behavior. Consequently, improving robustness only at the prompt or output level may be insufficient.

The brittleness of safety alignment is therefore a significant concern, particularly given the growing ecosystem of open-weight models, making stronger and more robust safety mechanisms for such models increasingly important.


Current defensive approaches do not directly address the sparsity of safety alignment and instead rely on methods such as external guardrails for user inputs or model responses~\cite{Inan2023llamaguard, robey2025smoothllm, kumar2024certifying, alon2023perplexity} to filter harmful prompts or responses before or after generation, or finetuning approaches that improve the general robustness of LLMs against jailbreak attacks using techniques such as negative preference optimization~\cite{zhao2025door} or chain-of-thought reasoning for safety alignment~\cite{hu2026alignmentweighted}. However, as our experiments show, these approaches still suffer from the same sparsity in safety alignment as the original models.

This work takes the first steps toward mitigating the sparsity of safety alignment in large language models. We propose \method (ABS-Align), which, at a high level, uses preference optimization to finetune a large language model toward more distributed safety alignment representations by temporarily and gradually deactivating the initially identified sparse safety-critical neurons. This process encourages a less sparse and concentrated distribution of safety-aligned behavior across the model, such that, during the later stages of finetuning, all originally identified safety neurons are temporarily deactivated while the model is optimized to maintain safety alignment through other neurons and representations. Rather than directly searching for new safety neurons, ABS-Align repeatedly removes the neurons with the strongest safety signal and finetunes the remaining model to re-introduce safety alignment. We use the term denser to describe safety-algined behavior that is less fragile and less dependent on a small set of safety-critical neurons that sparsity-based attacks detect, remove, or manipulate in order to jailbreak a target LLM.

Interestingly, our experiments show that the benefits of the less fragile safety-aligned model achieved by ABS-Align extend beyond the white-box attacks that motivate ABS-Align. Our experiments show that ABS-Align leads to broad improvements against a range of black-box attacks, despite receiving no data augmentation or explicit guidance for robustness against these attacks during finetuning. The finetuned model also exhibits greater honesty and truthfulness, while attacks that exploit the sparsity of safety alignment at both the neuron level and the activation and representation levels become significantly less effective. 

We perform extensive experiments across four models from three different vendors, specifically Google's Gemma, Meta's Llama model families, and Alibaba's Qwen. Our results show that ABS-Align consistently improves the models' safety alignment against a broad range of jailbreak attacks while having a negligible effect on their general capabilities and instruction-following ability.

The contributions of this work can be summarized as follows:

\begin{itemize}
    
    \item We propose ABS-Align, as a defense that mitigates the sparsity of safety alignment in LLMs by specifically finetuning the model to distribute its safety alignment across more parts of the network, improving robustness of safety alignment.
    
    \item By making safety-aligned LLMs less fragile, ABS-Align is the first to show increased robustness against not only pruning-based attacks, but also other forms of white-box attacks that conceptually rely on the sparsity of safety alignment. ABS-Align also generalizes to black-box attacks without any explicit data augmentation or guidance during finetuning.

    \item Our experiments provide causal evidence that finetuning with ABS-Align leads to a less sparse safety alignment and reduces the dependence of safety-aligned behavior on sparse model components. Specifically, we show that the finetuned models are substantially more robust to sparsity-based attacks that prune safety-critical neurons using two different jailbreak methods, TwinBreak and NeuroStrike. After fine-tuning, these methods' ability to jailbreak LLMs by pruning safety neurons is substantially degraded where even pruning 3 times more neurons compared to the undefended model shows minimal jailbreak success. Additionally, ablating a single refusal direction becomes substantially less effective, providing further evidence that the improved robustness extends beyond neuron-level interventions to attacks targeting the model's internal representations.


    \item We conduct a comprehensive set of experiments and evaluations across multiple models and dimensions, including attack success rates for white-box and black-box attacks, retention of general model capabilities across multiple benchmarks, and over-refusal rates. We compare our results with those of recent related work and show improvements across multiple evaluation dimensions, including reductions of approximately up to 90\% in attack success rate for some attacks.
\end{itemize}

Overall, ABS-Align provides a simple yet effective approach for mitigating fragility of safety alignment in large language models, improving robustness not only against attacks that exploit sparse safety-alignment representations but also against a broad range of prompt-level and black-box jailbreak attacks.



\textbf{Outline.} The remainder of the paper is structured as follows. Section~\ref{sec:background} provides the necessary background, and Section~\ref{sec:method} describes our method. Section~\ref{sec:experiments} presents the experimental evaluation and results, followed by a discussion in Section~\ref{sec:discussion}. Section~\ref{sec:related_work} reviews related work, and Section~\ref{sec:conclusion} concludes the paper.